% !TEX project = pdflatex
\documentclass[conference,a4paper]{IEEEtran}

\usepackage{packages}
\input{macros}

\begin{document}

\title{Computer Vision and Image Processing Project (2025/26)\\
       \Large Deep Learning for Traffic Sign Recognition using the GTSRB Dataset}

\author{

    \IEEEauthorblockN{Ana Oliveira}
    \IEEEauthorblockA{MEI, PG61510}
    \and
    \IEEEauthorblockN{André Miranda}
    \IEEEauthorblockA{MEI, PG60238}
    \and
    \IEEEauthorblockN{Maria Matos}
    \IEEEauthorblockA{MEI, PG61533}
}

\maketitle

\begin{abstract}
This report presents the first phase of a deep learning project focused on traffic sign recognition using the German Traffic Sign Recognition Benchmark (GTSRB) dataset. The main objective of this phase was to study the impact of different data augmentation techniques and image preprocessing strategies on the performance of convolutional neural networks (CNNs). Several augmentation configurations were evaluated, including geometric transformations, color-based augmentations, combined augmentations, and massive dataset expansion. Experimental results demonstrate that data augmentation significantly improves model generalization and robustness, achieving a maximum average test accuracy of approximately 98.38\%.
\end{abstract}

\begin{IEEEkeywords}
Deep Learning, Computer Vision, Traffic Sign Recognition, GTSRB, Convolutional Neural Networks, Data Augmentation, Image Processing, Ensemble Learning
\end{IEEEkeywords}

\section{Introduction}
Traffic sign recognition is an important task in computer vision and autonomous driving systems. Modern Advanced Driver Assistance Systems (ADAS) rely heavily on robust visual recognition models capable of detecting and classifying traffic signs under different environmental conditions such as illumination changes, rotations, occlusions, and motion blur.

The objective of this project is to explore deep learning techniques applied to the GTSRB dataset in order to achieve the highest possible classification accuracy. According to published benchmarks, the best reported accuracy on the dataset is approximately 99.82\%.

This first phase focuses mainly on:

\begin{itemize}
    \item Studying image preprocessing techniques;
    \item Applying static and dynamic data augmentation;
    \item Training multiple CNN models with different augmentation pipelines;
    \item Comparing the impact of augmentation strategies on model performance.
\end{itemize}

\section{Dataset}

The German Traffic Sign Recognition Benchmark (GTSRB) is a well-known benchmark dataset for traffic sign classification.

The dataset contains:
\begin{itemize}
    \item 43 traffic sign classes;
    \item More than 50,000 images;
    \item Images captured under real-world conditions;
    \item Significant variability in illumination, perspective, scale, and occlusion.
\end{itemize}

The dataset was divided into:
\begin{itemize}
    \item Training set;
    \item Validation set;
    \item Test set.
\end{itemize}

\subsection{Class Distribution}

One important characteristic of the dataset is the imbalance between classes. Some traffic signs contain considerably more samples than others, which can introduce bias during training.

The initial dataset analysis revealed:
\begin{itemize}
    \item Strong variability in the number of samples per class;
    \item Variability in image resolution;
    \item Presence of noisy samples and illumination changes.
\end{itemize}

\section{Methodology}

\subsection{Framework and Libraries}

The project was implemented using:
\begin{itemize}
    \item Python;
    \item PyTorch;
    \item Torchvision;
    \item NumPy;
    \item Matplotlib;
    \item Pandas.
\end{itemize}

\subsection{Image Preprocessing}

Before training, all images were resized to a fixed dimension in order to ensure consistency during batch processing.

The preprocessing pipeline included:
\begin{itemize}
    \item Image resizing;
    \item Conversion to tensor format;
    \item Normalization;
    \item Data type conversion to floating point.
\end{itemize}

Example preprocessing pipeline:

\begin{lstlisting}[language=Python]
transform_base = v2.Compose([
    v2.Resize((IMG_SIZE, IMG_SIZE)),
    v2.ToImage(),
    v2.ToDtype(torch.float32, scale=True)
])
\end{lstlisting}

\section{Model Architecture}

A convolutional neural network (CNN) architecture was implemented for traffic sign classification.

The network includes:
\begin{itemize}
    \item Convolutional layers;
    \item Batch normalization;
    \item ReLU activation functions;
    \item Max pooling layers;
    \item Fully connected layers;
    \item Dropout regularization.
\end{itemize}

The model was designed to balance:
\begin{itemize}
    \item Computational efficiency;
    \item Generalization capability;
    \item Classification accuracy.
\end{itemize}

\section{Data Augmentation Strategies \red{TODO}}

\section{Experimental Results \red{TODO} }

\section{Discussion}

The experiments demonstrate several important observations:

\begin{itemize}
    \item The baseline model already achieved strong performance due to the quality of the dataset and the CNN architecture.
    
    \item Geometric augmentation did not significantly improve performance, possibly because excessive transformations introduced unrealistic samples.
    
    \item Color augmentation slightly improved robustness while maintaining stable results.
    
    \item Combined augmentation introduced additional variability but may have increased training difficulty.
    
    \item Massive dataset expansion achieved the best overall performance, indicating that stronger regularization and diversity improved generalization.
\end{itemize}

The best average test accuracy obtained during Phase 1 was:

\[
98.38\%
\]

Although this result is below the state-of-the-art benchmark (99.82\%), it demonstrates that augmentation strategies positively impact model robustness and generalization.

\red{\section{Challenges and Limitations?}}

\section{Conclusion}

This first phase successfully explored the application of deep learning techniques to the GTSRB traffic sign dataset.

The experiments demonstrated that:
\begin{itemize}
    \item Data augmentation is essential for improving generalization;
    \item Different augmentation techniques have distinct impacts on performance;
    \item Stronger augmentation and dataset expansion yielded the best results.
\end{itemize}

The best model achieved approximately 98.38\% accuracy on the test set, providing a strong baseline for future improvements.

The second phase of the project will focus on ensemble learning techniques by combining the best-performing neural networks trained during Phase 1. The objective is to evaluate whether ensembles can improve classification accuracy, robustness, and generalization performance on the GTSRB dataset, potentially approaching state-of-the-art results.


\begin{thebibliography}{9}

\bibitem{gtsrb}
J. Stallkamp, M. Schlipsing, J. Salmen, and C. Igel,
\textit{The German Traffic Sign Recognition Benchmark: A multi-class classification competition},
International Joint Conference on Neural Networks (IJCNN), 2011.

\bibitem{pytorch}
PyTorch Documentation,
\url{https://pytorch.org/}

\bibitem{cnn}
Y. LeCun et al.,
\textit{Gradient-based learning applied to document recognition},
Proceedings of the IEEE, 1998.

\bibitem{augmentation}
A. Shorten and T. Khoshgoftaar,
\textit{A survey on Image Data Augmentation for Deep Learning},
Journal of Big Data, 2019.

\end{thebibliography}

\end{document}